\section{Cayley Transform}
\label{sec:cayley_transform}
Let $\tau \equiv i \frac{1-B}{1+B}$ transform as 
\begin{equation}
\tau \to \frac{a\tau+b}{c\tau+d}
\end{equation}
for
\begin{equation}
\begin{pmatrix}
a & b \\ c & d
\end{pmatrix}
\in \text{SL(2,R)}.
\end{equation}
Then the Cayley transform
\begin{equation}
B = \frac{i-\tau}{i+\tau}
\end{equation}
maps the upper half-plane parameterised by $\tau$ to the unit disk parameterised by $B$, and the SL(2,$\mathbb{R}$) transformation of $\tau$ to an SU(1,1) transformation of $B$.
More precisely,
\begin{equation}
B \to \frac{\alpha B + \beta}{\gamma B + \delta},
\end{equation}
where
\begin{equation}
\begin{pmatrix}
\alpha & \beta \\ \gamma & \delta
\end{pmatrix}
=
\begin{pmatrix}
(c-b)+i(a+d) & -(b-c)+i(d-a) \\ (b+c)+i(d-a) & (b-c)+i(a+d)
\end{pmatrix}
\end{equation}
is in SU(1,1) up to a sign.